\documentclass[12pt, a4paper]{article}

\usepackage[margin=2.5cm]{geometry}
\usepackage{hyperref}
\usepackage{xcolor}
\usepackage{listings}
\usepackage{titlesec}
\usepackage{enumitem}
\usepackage{fancyhdr}
\usepackage{mdframed}
\usepackage{parskip}
\usepackage{microtype}

\definecolor{primary}{HTML}{1D9E75}
\definecolor{secondary}{HTML}{534AB7}
\definecolor{codebg}{HTML}{F5F5F0}
\definecolor{codeborder}{HTML}{D3D1C7}
\definecolor{notebg}{HTML}{E1F5EE}
\definecolor{noteborder}{HTML}{0F6E56}
\definecolor{warnbg}{HTML}{FAEEDA}
\definecolor{warnborder}{HTML}{854F0B}
\definecolor{darktext}{HTML}{2C2C2A}
\definecolor{mutedtext}{HTML}{5F5E5A}
\definecolor{pykeyword}{HTML}{0F6E56}
\definecolor{pystring}{HTML}{BA7517}
\definecolor{pycomment}{HTML}{7F7D75}

\hypersetup{
  colorlinks=true,
  linkcolor=secondary,
  urlcolor=primary,
  citecolor=primary,
  pdftitle={Week 2 Report Agent Guide},
  pdfauthor={}
}

\titleformat{\section}{\Large\bfseries\color{darktext}}{}{0em}{}[\color{primary}\titlerule]
\titleformat{\subsection}{\large\bfseries\color{darktext}}{}{0em}{}
\titleformat{\subsubsection}{\normalsize\bfseries\color{mutedtext}}{}{0em}{}
\titlespacing{\section}{0pt}{2em}{0.8em}
\titlespacing{\subsection}{0pt}{1.4em}{0.4em}
\titlespacing{\subsubsection}{0pt}{1em}{0.3em}

\lstdefinestyle{plainstyle}{
  backgroundcolor=\color{codebg},
  basicstyle=\ttfamily\small\color{darktext},
  breaklines=true,
  frame=single,
  rulecolor=\color{codeborder},
  framesep=8pt,
  showstringspaces=false,
  tabsize=2
}

\lstdefinestyle{pythonstyle}{
  language=Python,
  backgroundcolor=\color{codebg},
  basicstyle=\ttfamily\small\color{darktext},
  keywordstyle=\color{pykeyword}\bfseries,
  stringstyle=\color{pystring},
  commentstyle=\color{pycomment}\itshape,
  breaklines=true,
  frame=single,
  rulecolor=\color{codeborder},
  framesep=8pt,
  showstringspaces=false,
  tabsize=4,
  numbers=left,
  numbersep=8pt
}

\lstdefinestyle{shellstyle}{
  backgroundcolor=\color{codebg},
  basicstyle=\ttfamily\small\color{darktext},
  breaklines=true,
  frame=single,
  rulecolor=\color{codeborder},
  framesep=8pt,
  showstringspaces=false,
  tabsize=2,
  morekeywords={python,pip,source}
}

\newmdenv[
  backgroundcolor=notebg,
  linecolor=noteborder,
  linewidth=3pt,
  topline=false,
  bottomline=false,
  rightline=false,
  innerleftmargin=12pt,
  innerrightmargin=12pt,
  innertopmargin=8pt,
  innerbottommargin=8pt
]{notebox}

\newmdenv[
  backgroundcolor=warnbg,
  linecolor=warnborder,
  linewidth=3pt,
  topline=false,
  bottomline=false,
  rightline=false,
  innerleftmargin=12pt,
  innerrightmargin=12pt,
  innertopmargin=8pt,
  innerbottommargin=8pt
]{warnbox}

\pagestyle{fancy}
\fancyhf{}
\fancyhead[L]{\small\color{mutedtext}Week 2 Report Agent}
\fancyfoot[C]{\small\color{mutedtext}\thepage}
\renewcommand{\headrulewidth}{0.4pt}
\renewcommand{\headrule}{\hbox to\headwidth{\color{codeborder}\leaders\hrule height \headrulewidth\hfill}}

\begin{document}

\begin{titlepage}
  \centering
  \vspace*{3cm}
  {\color{primary}\rule{\linewidth}{2pt}}
  \vspace{1cm}

  {\Huge\bfseries\color{darktext} Week 2: MCP-Enabled Report Agent}\\[0.5em]
  {\Large\color{mutedtext} Incremental guide from local files to MCP resources}

  \vspace{1cm}
  {\color{primary}\rule{\linewidth}{2pt}}

  \vspace{2cm}
  {\large\color{mutedtext}
    Extra concepts needed for \texttt{week-2/agent.py}\\
    and \texttt{week-2/mcp\_server.py}.
  }

  \vfill
  {\small\color{mutedtext} Version 1.0}
\end{titlepage}

\tableofcontents
\newpage

\section{How to Read This Document}

This document assumes you have already read Week 1. It does not repeat the introductory quarterly-report example, the use-case folder contract, or the full local pipeline walkthrough. Those are the foundation.

Week 2 is an incremental addendum. Read it as the answer to one question: \textit{what changes when the same use-case assets can be read through MCP resources instead of only through local file paths?}

\section{What Changes in Week 2}

Week 2 keeps the Week 1 pipeline model. The agent still resolves a use case, loads the task description and input data, discovers ordered steps, calls the model for each step, writes \texttt{context.json}, and saves the final output.

The new idea is that the agent can read use-case assets through two channels:

\begin{itemize}[leftmargin=1.5em]
  \item \textbf{\texttt{local}} -- read files directly from \texttt{use-cases/}, as in Week 1.
  \item \textbf{\texttt{mcp}} -- read the same assets through MCP resource URIs served by \texttt{week-2/mcp\_server.py}.
\end{itemize}

\begin{notebox}
\textbf{Week 2 is an extension, not a rewrite.} The prompt construction, step execution, context writing, and final report writing are the same idea as Week 1. The main change is how the agent fetches the files it needs.
\end{notebox}

\section{New Dependencies}

Week 1 needs LangChain and Ollama. Week 2 adds FastMCP:

\begin{lstlisting}[style=shellstyle]
pip install langchain-ollama fastmcp
\end{lstlisting}

If you use the repository's \texttt{requirements.txt}, \texttt{fastmcp} is already included.

FastMCP is used in two places:

\begin{itemize}[leftmargin=1.5em]
  \item \texttt{FastMCP} defines the MCP server and its resources in \texttt{mcp\_server.py}.
  \item \texttt{Client} lets \texttt{agent.py} read those resources when \texttt{AGENT\_TYPE = "mcp"}.
\end{itemize}

\section{MCP in This Project}

MCP stands for Model Context Protocol. In this project, you can think of it as a standard way to expose project assets to an agent. The MCP server presents files from \texttt{use-cases/} as named resources such as:

\begin{lstlisting}[style=plainstyle]
use-cases://
use-cases://quarterly-report/task
use-cases://quarterly-report/input
use-cases://quarterly-report/steps
use-cases://quarterly-report/step_01_data_structuring/spec
use-cases://quarterly-report/step_01_data_structuring/output
\end{lstlisting}

The Week 2 agent can ask the MCP client to read one of these resource URIs. The server maps the URI back to the appropriate file under \texttt{use-cases/}.

\begin{warnbox}
In this repository, MCP mode changes the read path only. Step outputs are still written locally to \texttt{context.json} and \texttt{final\_report.md} under the resolved use-case folder.
\end{warnbox}

\section{The MCP Server}

The MCP server is \texttt{week-2/mcp\_server.py}. It imports FastMCP, points at the same \texttt{use-cases/} directory, and registers resources.

\subsection{Server setup}

\begin{lstlisting}[style=pythonstyle]
import json
from pathlib import Path
from fastmcp import FastMCP

USE_CASES_DIR = Path(__file__).parent.parent / "use-cases"

mcp = FastMCP("report-agent")
\end{lstlisting}

\texttt{mcp = FastMCP("report-agent")} creates a named MCP server object. The resource decorators later in the file attach functions to this server.

\subsection{Path helpers}

\begin{lstlisting}[style=pythonstyle]
def _use_case_path(name: str) -> Path:
    path = USE_CASES_DIR / name
    if not path.exists():
        raise FileNotFoundError(f"Use case '{name}' not found")
    return path


def _step_path(name: str, step: str) -> Path:
    path = _use_case_path(name) / step
    if not path.exists():
        raise FileNotFoundError(f"Step '{step}' not found in use case '{name}'")
    return path
\end{lstlisting}

These helpers keep validation in one place. Resource functions call them before reading files.

\subsection{Listing use cases}

\begin{lstlisting}[style=pythonstyle]
@mcp.resource("use-cases://")
def list_use_cases() -> str:
    """List all available use case names."""
    names = sorted(p.name for p in USE_CASES_DIR.iterdir() if p.is_dir())
    return json.dumps(names, indent=2)
\end{lstlisting}

The root resource returns a JSON list of available use-case folders. The agent uses this during use-case resolution.

\subsection{Listing steps}

\begin{lstlisting}[style=pythonstyle]
@mcp.resource("use-cases://{name}/steps")
def list_steps(name: str) -> str:
    """List all step folder names for a given use case."""
    base = _use_case_path(name)
    steps = sorted(
        (p.name for p in base.iterdir() if p.is_dir() and p.name.startswith("step_")),
        key=lambda name: int(name.split("_")[1]),
    )
    return json.dumps(steps, indent=2)
\end{lstlisting}

This is the MCP version of Week 1's \texttt{fetch\_steps}. It returns JSON text because MCP resources return content, not Python objects. The agent parses the JSON after reading it.

\subsection{Reading task and input files}

\begin{lstlisting}[style=pythonstyle]
@mcp.resource("use-cases://{name}/task")
def get_task(name: str) -> str:
    """Return the agent task description for a use case."""
    return (_use_case_path(name) / "agent_task_description.md").read_text()


@mcp.resource("use-cases://{name}/input")
def get_input(name: str) -> str:
    """Return the input data for a use case."""
    return (_use_case_path(name) / "input_data.md").read_text()
\end{lstlisting}

These resources expose the two top-level files needed by every run.

\subsection{Reading step files}

\begin{lstlisting}[style=pythonstyle]
@mcp.resource("use-cases://{name}/{step}/spec")
def get_step_spec(name: str, step: str) -> str:
    """Return the spec for a specific step in a use case."""
    return (_step_path(name, step) / "spec.md").read_text()


@mcp.resource("use-cases://{name}/{step}/output")
def get_step_output(name: str, step: str) -> str:
    """Return the output format definition for a specific step in a use case."""
    return (_step_path(name, step) / "output.md").read_text()
\end{lstlisting}

These are the resources the agent reads before running each step.

\subsection{Optional prompt resource}

\begin{lstlisting}[style=pythonstyle]
@mcp.resource("use-cases://{name}/{step}/prompt")
def get_step_prompt(name: str, step: str) -> str:
    """Return the pre-rendered prompt for a step (available in Bedrock use cases)."""
    prompt_file = _step_path(name, step) / "prompt.md"
    if not prompt_file.exists():
        raise FileNotFoundError(f"No prompt.md found for step '{step}' in use case '{name}'")
    return prompt_file.read_text()
\end{lstlisting}

The current Week 2 agent does not use this resource during the local-or-MCP pipeline. It is present so use cases that contain pre-rendered \texttt{prompt.md} files can expose them through the same server.

\section{The Week 2 Agent}

\subsection{New imports and configuration}

Week 2 adds the FastMCP client and imports the server object:

\begin{lstlisting}[style=pythonstyle]
sys.path.insert(0, str(Path(__file__).parent))

from fastmcp import Client
from langchain_core.messages import HumanMessage, SystemMessage
from langchain_ollama import ChatOllama

from mcp_server import mcp

AGENT_TYPE = "local"  # either "local" or "mcp"
USE_CASES_DIR = Path(__file__).parent.parent / "use-cases"

_mcp_client: Client | None = None
\end{lstlisting}

\texttt{AGENT\_TYPE} chooses the read channel. Keep it as \texttt{"local"} to behave like Week 1. Set it to \texttt{"mcp"} to read through the MCP server.

\texttt{\_mcp\_client} stores the active client while an MCP run is in progress. This keeps the helper function signatures simple.

\subsection{Channel-aware asset fetching}

Week 1 had a local-only \texttt{fetch\_asset(path)}. Week 2 replaces it with a channel-aware version:

\begin{lstlisting}[style=pythonstyle]
async def fetch_asset(uri: str | Path, channel: str, type: str = "prompt") -> str:
    if channel == "local":
        path = Path(uri)
        return path.read_text() if path.exists() else ""

    if channel == "mcp":
        if _mcp_client is None:
            raise RuntimeError("MCP client is not initialized")
        contents = await _mcp_client.read_resource(str(uri))
        return contents[0].text

    raise ValueError(f"Unsupported channel '{channel}'")
\end{lstlisting}

The caller passes either a local path or an MCP resource URI. The rest of the pipeline does not need to know where the text came from.

\subsection{Channel-aware use-case and step discovery}

\begin{lstlisting}[style=pythonstyle]
async def fetch_use_cases(channel: str) -> list[str]:
    if channel == "local":
        return sorted(p.name for p in USE_CASES_DIR.iterdir() if p.is_dir())

    if channel == "mcp":
        return json.loads(await fetch_asset("use-cases://", channel))

    raise ValueError(f"Unsupported channel '{channel}'")
\end{lstlisting}

\begin{lstlisting}[style=pythonstyle]
async def fetch_steps(use_case: str, channel: str) -> list[str]:
    if channel == "local":
        use_case_dir = USE_CASES_DIR / use_case
        return sorted(
            (d.name for d in use_case_dir.iterdir() if d.is_dir() and d.name.startswith("step_")),
            key=lambda name: int(name.split("_")[1]),
        )

    if channel == "mcp":
        return json.loads(await fetch_asset(f"use-cases://{use_case}/steps", channel))

    raise ValueError(f"Unsupported channel '{channel}'")
\end{lstlisting}

These functions return Python lists in both modes. In local mode they build the list directly from the file system. In MCP mode they read JSON text from a resource and parse it.

\subsection{Mapping logical assets to locations}

The Week 2 agent adds \texttt{asset\_uri}. It translates logical asset names into either local paths or MCP URIs.

\begin{lstlisting}[style=pythonstyle]
def asset_uri(use_case: str, channel: str, asset: str, step_name: str | None = None) -> str | Path:
    if channel == "local":
        if asset == "task":
            return USE_CASES_DIR / use_case / "agent_task_description.md"
        if asset == "input":
            return USE_CASES_DIR / use_case / "input_data.md"
        if step_name is not None:
            return USE_CASES_DIR / use_case / step_name / f"{asset}.md"

    if channel == "mcp":
        if asset in {"task", "input"}:
            return f"use-cases://{use_case}/{asset}"
        if step_name is not None:
            return f"use-cases://{use_case}/{step_name}/{asset}"

    raise ValueError(f"Unsupported asset '{asset}' for channel '{channel}'")
\end{lstlisting}

This helper is the main design improvement over hard-coding paths everywhere. The pipeline can ask for \texttt{"task"}, \texttt{"input"}, \texttt{"spec"}, or \texttt{"output"} and let \texttt{asset\_uri} decide the concrete location.

\subsection{Running a step is almost unchanged}

The Week 2 \texttt{run\_step} receives one extra argument: \texttt{channel}. The only changed lines are the reads for \texttt{spec.md} and \texttt{output.md}.

\begin{lstlisting}[style=pythonstyle]
spec = await fetch_asset(asset_uri(use_case, channel, "spec", step_name), channel)
output_format = await fetch_asset(asset_uri(use_case, channel, "output", step_name), channel)
\end{lstlisting}

After those files are loaded, the same prompt-building and model invocation logic from Week 1 is used.

\section{MCP Client Lifecycle}

The agent should only create an MCP client when it actually needs MCP mode. That is handled by \texttt{run\_pipeline\_with\_client}.

\begin{lstlisting}[style=pythonstyle]
async def run_pipeline_with_client(use_case_input: str) -> str:
    global _mcp_client

    if AGENT_TYPE == "local":
        return await run_pipeline(use_case_input)

    async with Client(mcp) as client:
        _mcp_client = client
        try:
            return await run_pipeline(use_case_input)
        finally:
            _mcp_client = None
\end{lstlisting}

In local mode, no MCP client is created. In MCP mode, the client is opened once for the full pipeline run and then cleared in the \texttt{finally} block.

\begin{notebox}
\textbf{Why one client for the whole run?} The pipeline reads multiple resources: use cases, steps, task, input, specs, and output formats. Opening the client once keeps the lifetime clear and avoids repeated setup work.
\end{notebox}

\section{Running Week 2}

\subsection{Local mode}

With the default setting:

\begin{lstlisting}[style=pythonstyle]
AGENT_TYPE = "local"
\end{lstlisting}

run:

\begin{lstlisting}[style=shellstyle]
python week-2/agent.py "quarterly report"
\end{lstlisting}

The resolver maps \texttt{"quarterly report"} to the matching use-case folder, such as \texttt{quarterly-report}. This should behave like the Week 1 agent, while using the Week 2 code path.

\subsection{MCP mode}

Change the setting in \texttt{week-2/agent.py}:

\begin{lstlisting}[style=pythonstyle]
AGENT_TYPE = "mcp"
\end{lstlisting}

Then run the same command:

\begin{lstlisting}[style=shellstyle]
python week-2/agent.py "quarterly report"
\end{lstlisting}

The pipeline behavior should be the same from the user's point of view. Internally, the agent reads task assets through the MCP resources instead of direct file reads.


\section{Week 1 vs Week 2 Summary}

\begin{center}
\begin{tabular}{p{0.28\linewidth}p{0.30\linewidth}p{0.30\linewidth}}
\textbf{Topic} & \textbf{Week 1} & \textbf{Week 2} \\
\hline
Asset source & Local files only & Local files or MCP resources \\
Main script & \texttt{week-1/agent.py} & \texttt{week-2/agent.py} \\
Server file & None & \texttt{week-2/mcp\_server.py} \\
Fetcher & \texttt{fetch\_asset(path)} & \texttt{fetch\_asset(uri, channel)} \\
Mode switch & None & \texttt{AGENT\_TYPE} \\
New dependency & None beyond Week 1 & \texttt{fastmcp} \\
\end{tabular}
\end{center}

\section{Common Mistakes}

\begin{itemize}[leftmargin=1.5em]
  \item \textbf{Forgetting to install FastMCP.} MCP mode needs \texttt{fastmcp}; Week 1 does not.
  \item \textbf{Setting \texttt{AGENT\_TYPE} to an unsupported value.} It must be exactly \texttt{"local"} or \texttt{"mcp"}.
  \item \textbf{Breaking resource URI names.} The client and server must agree on URI patterns such as \texttt{use-cases://\{name\}/input}.
  \item \textbf{Assuming MCP mode writes remotely.} In this project, writes still go to the local use-case folder.
  \item \textbf{Changing the use-case folder format.} Both channels still depend on \\ \texttt{agent\_task\_description.md}, \texttt{input\_data.md}, and numbered step folders.
\end{itemize}

\end{document}
